% 数学ノート集・共通テンプレート v2（積分論2026年度第1部の版面を参考）
% A5判。章内で定義・定理・例・注意等の番号を共通化。
% コンパイル： platex -interaction=nonstopmode math_note_template_v2.tex (2回)
%              dvipdfmx math_note_template_v2.dvi
\documentclass[a5paper,10pt,onecolumn]{jsarticle}
\usepackage[dvipdfmx,top=22mm,bottom=20mm,inner=19mm,outer=18mm,headheight=13pt,headsep=6mm]{geometry}
\usepackage{amsmath,amssymb,amsthm,mathtools}
\usepackage{enumitem,fancyhdr,titlesec}
\usepackage[dvipdfmx,bookmarks=true,colorlinks=false]{hyperref}
\usepackage{pxjahyper}

% ----- 差し替える文書情報 -----
\newcommand{\NoteTitle}{離散時間金融市場と無裁定価格理論}
\newcommand{\NoteSeries}{数学ノート集}

% ----- 文字・組版 -----
% 参考PDFの質素な明朝＋Computer Modernの印象を引き継ぎながら，
% A5判で窮屈にならないよう本文の行間を少し確保する．
\linespread{1.13}
\setlength{\parindent}{1zw}
\setlength{\parskip}{0pt}
\setlength{\abovedisplayskip}{6pt plus 2pt minus 2pt}
\setlength{\belowdisplayskip}{6pt plus 2pt minus 2pt}
\setlength{\abovedisplayshortskip}{3pt plus 2pt minus 1pt}
\setlength{\belowdisplayshortskip}{4pt plus 2pt minus 1pt}
\setlist[enumerate]{topsep=3pt,itemsep=1pt,parsep=0pt,leftmargin=2.4zw}
\titleformat{\section}{\normalfont\normalsize\bfseries}{\thesection}{0.8zw}{}
\titlespacing*{\section}{0pt}{2.0ex plus .4ex}{.9ex}
\titleformat{\subsection}{\normalfont\normalsize\bfseries}{\thesubsection}{0.7zw}{}
\titlespacing*{\subsection}{0pt}{1.5ex plus .3ex}{.6ex}
\raggedbottom

% ----- 柱とページ番号：ページ上部、罫線なし -----
\fancypagestyle{mathnote}{%
  \fancyhf{}%
  \fancyhead[L]{\footnotesize \NoteSeries\ （\NoteTitle）}%
  \fancyhead[R]{\footnotesize \thepage}%
  \renewcommand{\headrulewidth}{0pt}%
  \renewcommand{\footrulewidth}{0pt}%
}
\pagestyle{mathnote}

% ----- 節内通し番号の定義・定理群 -----
\newtheoremstyle{mathnotetheorem}
 {4.5pt}{4.5pt}{\normalfont}{}%
 {\bfseries}{}{0.65zw}%
 {\thmname{#1}\thmnumber{ #2}\thmnote{（#3）}}
\theoremstyle{mathnotetheorem}
\newtheorem{theorem}{定理}[section]
\newtheorem{definition}[theorem]{定義}
\newtheorem{lemma}[theorem]{補題}
\newtheorem{proposition}[theorem]{命題}
\newtheorem{corollary}[theorem]{系}
\newtheorem{example}[theorem]{例}
\newtheorem{remark}[theorem]{注意}
\renewcommand{\proofname}{証明}
\renewcommand{\qedsymbol}{$\square$}
\makeatletter
\renewenvironment{proof}[1][\proofname]{%
  \par\pushQED{\qed}\normalfont\topsep4\p@\@plus2\p@\relax
  \trivlist\item[\hskip\labelsep\bfseries #1]\ignorespaces
}{\popQED\endtrivlist\@endpefalse}
\makeatother

% ----- 使うノートで適宜拡張する記号 -----
\newcommand{\R}{\mathbb{R}}
\newcommand{\norm}[1]{\lVert #1\rVert}
\newcommand{\abs}[1]{\lvert #1\rvert}

\begin{document}
\section*{到達目標と前提}
有限状態モデルの無裁定条件を確かめ，
一期間二項市場で条件付き請求権を株式と銀行口座で
実際に複製し，その価格を求められるようになる．
前提は線形代数，有限確率空間，期待値である．


\section{有限状態の金融市場}

本稿では $\Omega$ を有限集合とし，
確率測度 $P$ は全ての $\omega\in\Omega$ に
$P(\{\omega\})>0$ を割り当てると仮定する．
時刻は $t=0,1,\dots,T$ とし，
$\mathcal F_0\subset\cdots\subset\mathcal F_T
=\mathcal P(\Omega)$ をフィルトレーションとする．
各 $\mathcal F_t$ の原子を時刻 $t$ の節点という．

\begin{definition}
銀行口座の価値を $B_t>0$，
$d$ 個の危険資産の価値を
$S_t=(S_t^1,\dots,S_t^d)$ とする．
各 $B_t,S_t$ は $\mathcal F_t$-可測とする．
割引価格を
\[
 \widetilde S_t^i=\frac{S_t^i}{B_t},
 \qquad \widetilde S_t=(\widetilde S_t^1,\dots,\widetilde S_t^d)
\]
と定める．
$B_t$ は正の確定数であると仮定する．
\end{definition}

\begin{definition}
戦略とは，各 $t=1,\dots,T$ に対し
$\mathcal F_{t-1}$-可測な株式保有量
$\theta_t\in\mathbb R^d$ と，
対応する銀行口座保有量を定めることをいう．
戦略が\emph{自己金融的}であるとは，
再配分によって外部から資金を投入・引き出ししないことをいう．
割引資産価値を $\widetilde V_t=V_t/B_t$ と書く．
\end{definition}

\begin{proposition}[自己金融条件]\label{prop:self-financing}
自己金融戦略の割引資産価値は
\[
 \widetilde V_t
 =\widetilde V_0+
 \sum_{k=1}^t\theta_k\cdot
 (\widetilde S_k-\widetilde S_{k-1})
 \qquad(t=1,\dots,T)
 \tag{SF}
\]
を満たす．
逆に初期資本 $\widetilde V_0$ と任意の予測可能な
株式保有量 $(\theta_k)$ に対し，
銀行口座の保有量を適切に選べば
式 (SF) を満たす自己金融戦略が存在する．
\end{proposition}
\begin{proof}
時刻 $k-1$ の再配分直後の銀行口座保有量を
$\eta_k$ とすると，
時刻 $k-1$ での割引価値は
$\eta_k+\theta_k\cdot\widetilde S_{k-1}$ である．
時刻 $k$ まで保有量を変更しなければ
割引価値は
$\eta_k+\theta_k\cdot\widetilde S_k$ になる．
自己金融条件から前者は
$\widetilde V_{k-1}$ に等しく，
後者は $\widetilde V_k$ に等しいから
\[
 \widetilde V_k-\widetilde V_{k-1}
 =\theta_k\cdot(\widetilde S_k-\widetilde S_{k-1}).
\]
これを足し合わせて (SF) を得る．
逆に $(\theta_k)$ と $\widetilde V_0$ が与えられたら，
式 (SF) によって割引価値を帰納的に定め，
$\eta_k=\widetilde V_{k-1}
 -\theta_k\cdot\widetilde S_{k-1}$ と置けば，
各時点の再配分前後の価値は一致する．
$\eta_k$ は $\mathcal F_{k-1}$-可測であるから
予測可能な自己金融戦略となる．
\end{proof}

\begin{definition}
初期資本零の自己金融戦略で，
\[
 \widetilde V_T\ge0\ \text{（全ての状態で）},\qquad
 P(\widetilde V_T>0)>0
\]
を満たすものを\emph{裁定取引}という．
裁定取引が存在しないことを\emph{無裁定}（NA）という．
\end{definition}

\section{一期間市場の無裁定基本定理}

一期間の有限市場を考える．
状態を $\omega_1,\dots,\omega_m$ とし，
割引資産価格の変化を
\[
 v_i=\widetilde S_1(\omega_i)-\widetilde S_0
 \in\mathbb R^d
\]
と書く．
$\mathcal F_0$ は自明な $\sigma$-代数とする．

\begin{lemma}[正確な有限次元の分離]\label{lem:finite-alternative}
$v_1,\dots,v_m\in\mathbb R^d$ に対し，
次は同値である．
\begin{enumerate}
\item $\theta\in\mathbb R^d$ が
$\theta\cdot v_i\ge0$ を全ての $i$ で満たすなら，
$\theta\cdot v_i=0$ が全ての $i$ で成り立つ．
\item 正の実数 $q_1,\dots,q_m>0$ が存在して
$\sum_iq_i=1$ および
$\sum_iq_iv_i=0$ を満たす．
\end{enumerate}
\end{lemma}
\begin{proof}
(2) を仮定する．
$\theta\cdot v_i\ge0$ なら
\[
 0=\theta\cdot\sum_iq_iv_i
 =\sum_iq_i(\theta\cdot v_i)
\]
である．各 $q_i>0$ なので
$\theta\cdot v_i=0$ が全ての $i$ で成立する．

逆に (1) を仮定する．
\[
 L=\{(\theta\cdot v_1,\dots,\theta\cdot v_m):
 \theta\in\mathbb R^d\}\subset\mathbb R^m
\]
と置く．
$L$ は有限次元線形部分空間であり閉である．
仮定は $L\cap[0,\infty)^m=\{0\}$ と同値である．
標準単体
\[
 \Delta=\{u\in[0,\infty)^m:\sum_iu_i=1\}
\]
はコンパクト凸集合で，$L\cap\Delta=\varnothing$ である．
従って $\Delta$ と $L$ の距離は正である．
$\Delta$ 上で
$\norm{u-P_Lu}$ を最小にする元 $u_*$ を取り，
$w=u_*-P_Lu_*\in L^\perp$ と置く．
$w\ne0$ である．
任意の $u\in\Delta$ と $0<t<1$ について
$u_*+t(u-u_*)\in\Delta$ であり，
$u_*$ の最小性から
\[
 \norm{P_{L^\perp}u_*}^2
 \le\norm{P_{L^\perp}(u_*+t(u-u_*))}^2.
\]
展開して $t\downarrow0$ とすると
\[
 \langle w,u-u_*\rangle\ge0.
\]
また $\langle w,u_*\rangle=\norm w^2>0$ なので
$\langle w,u\rangle\ge\norm w^2$ が全ての $u\in\Delta$ で成り立つ．
$\Delta$ の各頂点である標準基底ベクトルを代入すると
$w_i\ge\norm w^2>0$ となる．
$w\in L^\perp$ だから
$\sum_iw_iv_i=0$ である．
$q_i=w_i/\sum_jw_j$ と置けば
(2) の条件を満たす．
\end{proof}

\begin{definition}
$Q$ が $P$ と\emph{同値}であるとは，
$P(A)=0$ と $Q(A)=0$ が同値であることをいう．
有限状態で $P(\omega_i)>0$ なら，
$Q$ が同値であることは $Q(\omega_i)>0$ が
全ての $i$ で成り立つことと同値である．
\end{definition}

\begin{theorem}[一期間市場の無裁定基本定理]
\label{thm:one-period-ftap}
一期間の有限市場が無裁定であるための必要十分条件は，
$P$ と同値な確率測度 $Q$ が存在して
\[
 E_Q[\widetilde S_1]=\widetilde S_0
\]
となることである．
\end{theorem}
\begin{proof}
無裁定の定義は
$\theta\cdot v_i\ge0$ が全ての $i$ で成り立ち，
少なくとも一つの $i$ で真に正となる
$\theta$ が存在しないことである．
これは補題\ref{lem:finite-alternative}の (1) に等しい．
従って正の重み $q_i$ で
$\sum_iq_i=1$，
$\sum_iq_iv_i=0$ となるものが存在することと同値である．
$Q(\omega_i)=q_i$ と定めれば $Q$ は $P$ と同値であり，
$\sum_iq_iv_i=0$ は
$E_Q[\widetilde S_1]=\widetilde S_0$ と同値である．
\end{proof}

\section{多期間市場とマルチンゲール測度}

\begin{definition}
$Q$ が $P$ と同値で，
割引価格過程 $(\widetilde S_t)$ が
$Q$-マルチンゲール，すなわち
\[
 E_Q[\widetilde S_t\mid\mathcal F_{t-1}]
 =\widetilde S_{t-1}\qquad(t=1,\dots,T)
\]
を満たすとき，
$Q$ を\emph{同値マルチンゲール測度}（EMM）という．
\end{definition}

\begin{proposition}[マルチンゲール測度と無裁定]
\label{prop:emm-no-arb}
有限多期間市場にEMM $Q$ が存在すれば，
その市場は無裁定である．
\end{proposition}
\begin{proof}
初期資本零の自己金融戦略を取る．
式 (SF) より
\[
 \widetilde V_T
 =\sum_{k=1}^T\theta_k\cdot
  (\widetilde S_k-\widetilde S_{k-1}).
\]
$\theta_k$ は $\mathcal F_{k-1}$-可測であり，
状態数が有限なので全ての項は可積分である．
条件付き期待値の塔の性質から
\begin{align*}
 E_Q[\theta_k\cdot(\widetilde S_k-\widetilde S_{k-1})]
 &=E_Q[\theta_k\cdot
   E_Q[\widetilde S_k-\widetilde S_{k-1}
       \mid\mathcal F_{k-1}]]\\
 &=0.
\end{align*}
従って $E_Q[\widetilde V_T]=0$ である．
もし $\widetilde V_T\ge0$ かつ
$P(\widetilde V_T>0)>0$ なら，
$Q\sim P$ により $Q(\widetilde V_T>0)>0$
となり，非負確率変数の期待値が正となる．
これは矛盾する．
\end{proof}

\begin{lemma}[節点ごとの無裁定]\label{lem:local-na}
有限多期間市場が無裁定なら，
任意の時刻 $t<T$ と $\mathcal F_t$ の原子 $A$ に対し，
$A$ から次の時刻へ進む一期間の条件付き市場は
無裁定である．
\end{lemma}
\begin{proof}
$A$ における条件付き市場に裁定ポジション
$\theta\in\mathbb R^d$ があると仮定する．
$t+1$ の時刻だけ $A$ 上で $\theta$ を保有し，
それ以外の時刻と節点では零を保有する戦略を考える．
$\theta_{t+1}=\theta\,1_A$ は
$\mathcal F_t$-可測なので予測可能である．
初期資本零に対する式 (SF) により，
この戦略の最終割引利益は
\[
 1_A\,\theta\cdot
 (\widetilde S_{t+1}-\widetilde S_t)
\]
である．
$A$ の内部で非負かつ正の値を取る状態があり，
$A$ の外では零だから全体の裁定となる．
仮定に反するので条件付き市場は無裁定である．
\end{proof}

\begin{theorem}[有限多期間市場の無裁定基本定理]
\label{thm:finite-ftap}
有限状態・有限期間市場について，
無裁定であることとEMMが存在することは同値である．
\end{theorem}
\begin{proof}
EMMの存在から無裁定は
命題\ref{prop:emm-no-arb}による．
逆に無裁定とする．
補題\ref{lem:local-na}によって任意の時刻
$t<T$ と $\mathcal F_t$ の原子 $A$ 上の
条件付き一期間市場は無裁定である．
その次の時刻の子節点を
$A_1,\dots,A_r$ とし，
割引価格の変化を $v_i$ とする．
一期間の定理\ref{thm:one-period-ftap}から，
全て正の重み $q(A_i\mid A)>0$ で
\[
 \sum_iq(A_i\mid A)=1,\qquad
 \sum_iq(A_i\mid A)v_i=0
\]
となるものが存在する．
有限個の節点全てについてこのような重みを選ぶ．
$\mathcal F_0$ の原子が複数ある場合は，
それぞれに正の初期確率を割り当てる．
次に各時刻で親節点から子節点へ
$q(A_i\mid A)$ を掛け合わせることにより，
終端状態の確率 $Q(\omega)>0$ を定める．
各節点からの遷移重みの和は1なので，
これは全確率が1となる確率測度である．
全状態に正の重みがあるため $Q\sim P$ である．
さらに各節点 $A$ で
\[
 E_Q[\widetilde S_{t+1}-\widetilde S_t
       \mid\mathcal F_t](A)
 =\sum_iq(A_i\mid A)v_i=0.
\]
従って $\widetilde S$ は $Q$-マルチンゲールであり，
$Q$ はEMMである．
\end{proof}

\section{複製価格とマルチンゲール}

\begin{definition}
満期 $T$ の非負とは限らない確率変数 $H$ を
\emph{請求権}と呼ぶ．
自己金融戦略の満期価値が $H$ に等しくなるとき，
その戦略は $H$ を\emph{複製}するという．
\end{definition}

\begin{theorem}[複製可能な請求権の価格]\label{thm:replication}
有限多期間市場にEMM $Q$ が存在し，
請求権 $H$ が初期資本 $V_0$ の自己金融戦略で
複製可能であるとする．
このとき
\[
 \frac{V_0}{B_0}=E_Q\left[\frac H{B_T}\right].
\]
特に複製可能な請求権の初期価格は，
EMMの選択によらず一意に定まる．
\end{theorem}
\begin{proof}
自己金融式 (SF) の両辺について
$Q$-期待値を取る．
命題\ref{prop:emm-no-arb}の証明と同様に
各割引利益の期待値は零なので
\[
 E_Q[\widetilde V_T]=\widetilde V_0=\frac{V_0}{B_0}
\]
となる．
複製戦略の満期価値は $V_T=H$ であるから
$\widetilde V_T=H/B_T$ となり表示式を得る．
$Q$ は任意のEMMだったので，
右辺はEMMによらず同じ値となる．
\end{proof}

\section{完全市場と二項モデル}

\begin{definition}
全ての請求権 $H:\Omega\to\mathbb R$ が
自己金融戦略で複製可能である市場を\emph{完備市場}という．
\end{definition}

\begin{theorem}[一期間市場の完備性]\label{thm:complete-oneperiod}
無裁定な一期間有限市場では，
次の条件は同値である．
\begin{enumerate}
\item 市場は完備である．
\item EMMは一意である．
\item $\mathbb R^m$ は
\[
 (1,\dots,1),\quad
 (v_1^{(j)},\dots,v_m^{(j)})\quad(j=1,\dots,d)
\]
によって張られる．
ここで $v_i^{(j)}$ は状態 $\omega_i$ での
割引価格変化 $v_i$ の第 $j$ 成分である．
\end{enumerate}
\end{theorem}
\begin{proof}
初期割引資本 $c$ と保有量 $\theta$ による
満期割引価値は
$(c+\theta\cdot v_i)_{i=1}^m\in\mathbb R^m$
である．
全ての請求権が複製可能であることは，
これらの価値ベクトルが $\mathbb R^m$ を尽くすこと，
すなわち (3) と同値である．

EMMの確率ベクトル $q=(q_i)$ は
\[
 q_i>0,\qquad\sum_iq_i=1,\qquad
 \sum_iq_iv_i=0
\]
を満たす．
そのような正の解 $q_*$ が少なくとも一つ存在するのは
無裁定基本定理による．
上の等式の係数行列の各行が
(3) に挙げたベクトルなので，
その張る空間が $\mathbb R^m$ なら
線形方程式の解は高々一つであり，
EMMは一意である．
逆に張る空間が真部分空間なら，
それに直交する非零ベクトル $h\in\mathbb R^m$ が存在する．
十分小さい $\varepsilon>0$ に対して
$q_\pm=q_*\pm\varepsilon h$ はどちらも
全ての成分が正であり，
\[
 \sum_i(q_\pm)_i=1,\qquad
 \sum_i(q_\pm)_iv_i=0
\]
を満たす．
二つは相異なるEMMなので一意性に反する．
従って (2) から (3) が従う．
\end{proof}

\begin{example}[一期間二項モデル]
銀行口座は $B_0=1$，$B_1=R>0$，
株価は $S_0=s>0$ とする．
$S_1=su$ または $S_1=sd$ となり，
$u>d>0$ とする．
この市場が無裁定となる条件は
\[
 d<R<u
\]
である．
このときEMMは一意であり，上昇確率は
\[
 q=\frac{R-d}{u-d}
\]
となる．
\end{example}
\begin{proof}
割引株価の上昇・下降時の変化はそれぞれ
$s(u/R-1)$ と $s(d/R-1)$ である．
無裁定はこれらが異符号であることに等しく，
$u>R>d$ となる．
一期間のEMM条件は
\[
 q\frac{su}{R}+(1-q)\frac{sd}{R}=s
\]
である．
これを解くと $q=(R-d)/(u-d)$ となり，
無裁定の条件から $0<q<1$ である．
二つの状態に対し現金と株式の二つの
独立な満期価値ベクトルがあるため，
定理\ref{thm:complete-oneperiod}から市場は完備である．
\end{proof}

\begin{remark}
有限状態の無裁定基本定理は，
状態ごとの割引利益の線形空間と非負錐の分離によって
証明できる．
連続時間のBlack--Scholesモデルでは，
これに対応する役割をGirsanovの定理と
マルチンゲール表現定理が担う．
\end{remark}

\section{一期間市場で複製戦略を構成する}
\begin{example}[株式オプションの無裁定価格]
時刻0の株価を $S_0=100$ とし，
時刻1には $S_u=120$ または $S_d=80$
となる一期間二項市場を考える．
銀行口座は時刻0の1単位が時刻1に $R=1.05$
となる．$80<RS_0=105<120$ なので無裁定である．
行使価格 $K=100$ のヨーロピアン・コールの満期支払は
\[
 H_u=(120-100)^+=20,\quad H_d=(80-100)^+=0.
\]
時刻0に株式を $\Delta$ 単位，銀行口座へ $B$ の金額を
置く戦略で複製するには
\[
 120\Delta+1.05 B=20,\qquad
 80\Delta+1.05 B=0
\]
が必要十分である．差を取ると
$\Delta=20/(120-80)=1/2$．
下の式に戻すと $B=-40/1.05=-800/21$．
したがって無裁定価格は
\[
 V_0=100\Delta+B
 =50-\frac{800}{21}=\frac{250}{21}.
\]
\end{example}
\begin{proposition}[リスク中立確率での検算]
上の価格は，リスク中立確率
$q=(105-80)/(120-80)=5/8$ を用いた
割引期待値 $R^{-1}E_QH$ と一致する．
\end{proposition}
\begin{proof}
$0<q<1$ であり
$q\,120+(1-q)\,80=105=RS_0$ だから
割引株価は $Q$ の下でマルチンゲールである．
\[
 R^{-1}E_QH
 =\frac1{1.05}\left(\frac58\cdot20
                  +\frac38\cdot0\right)
 =\frac{250}{21}.
\]
この値が複製費用に等しいことを確認した．
\end{proof}
\section*{到達確認問題と解答}
\noindent\textbf{問題．}
同じ市場で満期支払が上状態10，下状態30となる
請求権を複製する株式保有量と銀行口座への投資額を求めよ．

\medskip\noindent\textbf{解答．}
$120\Delta+1.05B=10$ と
$80\Delta+1.05B=30$ の差から
$\Delta=-1/2$．第二式より
$1.05B=70$ なので $B=200/3$．
初期費用は $100(-1/2)+200/3=50/3$．
リスク中立期待値からも
$R^{-1}((5/8)10+(3/8)30)
 =(20/21)(35/2)=50/3$ と一致する．

\end{document}
